An LLM-guided refinement loop for differential testing of Dart JSON
deserializers found divergence in 2.9\% of generated documents on its own,
but in 28.8\% when a human gave it four observations from a manual
characterization: the bottleneck was domain knowledge, and it came from a
person. We ask whether the fuzzer can acquire that knowledge itself. Holding
the loop fixed and varying only a 1{,}500-character knowledge slot, we
compare sources in a pre-registered study: none, the human's observations,
a summary written by an LLM agent that reads the deserializers' code, and a
summary written by an agent that designs and runs 20 black-box probe
documents. On Dart ($n{=}5$ seeds per arm), human knowledge reaches 30.9\%,
code-read knowledge 11.6\% and no knowledge 1.5\% (each pair
$p \approx 0.008$, Holm-adjusted $p \approx 0.048$). Unguided probing
reaches 10.5\%, not detectably above no knowledge ($p \approx 0.69$): its
probes never remove a required top-level field or leave the
64-bit integer range, and it
succeeds on one seed in five, by accident. A pre-registered extension tests
two remedies. Requiring the agent to cover generic test categories first
raises probing to 20.2\% ($p \approx 0.032$ against no knowledge;
$0.16$ after correction across six tests), with the missing-field pattern
found in every run; a larger model raises it to 12.1\%. Neither reaches
the human arm's mean, and both still leave the integer-overflow pattern mostly
unfound, because the probes never reach it. Across all arms, which known
divergence patterns the loop finds tracks which facts the knowledge states.
On Kotlin (Gson, Moshi, kotlinx.serialization, Jackson), where divergence is
common, probing helps without guidance (67.1\% against 50.0\% over ten
seeds, $p \approx 0.007$) and more with it (81.6\%). Self-characterization
works where interesting behaviour is dense, and needs systematic coverage
where it is narrow.
